\chapter{مرور ادبیات‌}\label{ch2}
الگوریتم‌های مسیریابی بر اساس میزان شناخت از محیط به دو دسته سراسری\LTRfootnote{\lr{global} }  (محیط ایستا/شناخته‌شده) و محلی\LTRfootnote{\lr{local} }  (محیط پویا/ناشناخته) تقسیم می‌شوند. همچنین می‌توان آن‌ها را به الگوریتم‌های جستجوی کلاسیک و الگوریتم‌های اکتشافی دسته‌بندی کرد. الگوریتم \lr{A*} یکی از پایه‌ای‌ترین روش‌ها برای ارزیابی کوتاه‌ترین مسیر در محیط‌های ایستا است \cite{karur2021survey, dudi2020shortest}. در حالی که الگوریتم‌های پویا برای چنین محیط‌هایی کارایی ندارند.

الگوریتم دیجسترا\LTRfootnote{\lr{Dijkstra} }  اگرچه کوتاه‌ترین مسیر را می‌یابد، اما به دلیل جستجوی کورکورانه و پیمایش کامل نقشه، محاسبات سنگین و کارایی پایینی دارد \cite{karur2021survey}. الگوریتم \lr{A} و نسخه‌های مختلف آن به عنوان الگوریتم‌های پیشرفته برای غلبه بر این محاسبات سنگین در محیط‌های ایستا معرفی شدند \cite{liu2017global}. با این حال، \lr{A*} نیز به دلیل نیاز به بررسی گره‌های متعدد، زمان محاسبه بالایی دارد که با بزرگ‌شدن مقیاس نقشه، کارایی آن کاهش می‌یابد \cite{wang2021overview}. علاوه بر این، استفاده از نقشه‌های شبکه‌ای\LTRfootnote{\lr{rasterizes maps} }  در این الگوریتم منجر به تولید مسیرهای ناصاف با چرخش‌های قائم‌الزاویه متعدد می‌شود که قابلیت اطمینان آن را پایین می‌آورد. این الگوریتم برای محاسبه هزینه مسیر از توابع اکتشافی نظیر فاصله منهتنی \cite{singh2018constrained}، فاصله اقلیدسی و فاصله قطری \cite{yang2015path} بهره می‌برد.

محققان روش‌های مختلفی برای بهبود کارایی \lr{A*} ارائه داده‌اند؛ از جمله: وزن‌دهی به توابع ارزیابی و علامت‌گذاری گره‌های مجاور موانع به عنوان «نقاط نامطمئن» \cite{qingji2005feasible}، بهینه‌سازی نحوه ذخیره‌سازی گره‌ها در آرایه‌ها برای جستجوی سریع‌تر  \cite{peng2015robot}و محاسبه زودهنگام توابع اکتشافی برایماهی زمان اجرا \cite{guruji2016time}؛   و هیچ‌کدام به بهبود استحکام\LTRfootnote{\lr{robustness } }  \lr{A*} الگوریتم نپرداخته‌اند.

برای جلوگیری از برخورد، مفهوم «فاصله ایمن» با ایجاد یک مرز دایره‌ای حول وسایل نقلیه سطحی پیشنهاد شده است. هر چند این روش موانع نقشه را به خوبی لحاظ نمی‌کند. برای رفع مشکل ناصافی مسیرها در \lr{A*}، روش‌هایی مانند ادغام مسیریابی سراسری با ویژگی‌های محیط محلی و حذف گره‌های اضافی ارائه شده است \cite{zhang2018dynamic}. همچنین استفاده از تکنیک‌های هموارسازی مانند درون‌یابی اسپالین مکعبی\LTRfootnote{\lr{Cubic Spline Interpolation}}  \cite{song2019smoothed}، برای کاهش چرخش‌های تند و خطر برخورد بسیار مؤثر بوده‌اند. علاوه بر این، نسخه‌های اصلاح‌شده‌ای مانند \lr{Basic Theta} و \lr{JPS} \cite{hart1968formal} و همچنین \lr{A*} برای وسایل نقلیه خودران توسعه یافته‌اند که امکان یافتن مسیر بهینه را فراهم می‌کنند.

روش‌های فعلی بهبودیافته از الگوریتم \lr{A*}، تنها روی یکی از دو فاکتور «کارایی» یا «استحکام» تمرکز کرده‌اند. با توجه به کاربرد گستره‌های متحرک در دنیای واقعی، جای خالی الگوریتمی جامع که همزمان دارای استحکام قوی و کارایی بالا باشد، احساس می‌شود؛ موضوعی که پتانسیل کاربردی و ارزش تجاری بالایی در بخش صنعت دارد.
